% Part: incompleteness
% Chapter: incompleteness-provability
% Section: introduction

\documentclass[../../../include/open-logic-section]{subfiles}

\begin{document}

\olfileid[pt]{inc}{inp}{int}

\olsection{Introdução}

Hilbert pensava que um sistema de axiomas para uma estrutura matemática,
como os números naturais, seria inadequado a menos que permitisse derivar
todas as afirmações verdadeiras sobre a estrutura. Combinada com seu
interesse posterior por sistemas formais de dedução, essa ideia sugere que
ele julgava necessário garantir que, digamos, os sistemas formais usados
para raciocinar sobre os números naturais fossem não apenas consistentes,
mas também \emph{completos}, isto é, que toda afirmação em sua linguagem ou
ela própria ou sua negação fosse !!{derivable}. O primeiro teorema da
incompletude de G\"odel mostra que não existe tal sistema de axiomas: não há
sistema formal completo, consistente e !!{axiomatizable} para a aritmética.
Na verdade, nenhuma teoria matemática ``suficientemente forte'',
consistente e !!{axiomatizable} é completa.

Um objetivo ainda mais importante de Hilbert, peça central de seu programa
para justificar a matemática moderna (``clássica''), era encontrar provas
finitárias de consistência para sistemas formais que representassem o
raciocínio clássico. No que diz respeito ao programa de Hilbert, portanto,
o segundo teorema da incompletude de G\"odel foi um golpe muito maior. O
segundo teorema da incompletude pode ser enunciado em termos vagos, assim
como o primeiro. Em linhas gerais, ele afirma que nenhuma teoria da
aritmética suficientemente forte pode provar a própria consistência.
Teremos de entender ``suficientemente forte'' como algo que inclua um pouco
mais que~$\Th{Q}$.

A ideia por trás da prova original de G\"odel do teorema da incompletude
pode ser encontrada no paradoxo de Epimênides. Epimênides, um cretense,
afirmou que todos os cretenses são mentirosos; uma forma mais direta do
paradoxo é a afirmação ``esta sentença é falsa''. Em essência, substituindo
verdade por !!{derivability}, G\"odel conseguiu formalizar !!a{sentence}
que, de maneira indireta, afirma não ser ela própria !!{derivable}. Se essa
!!{sentence} fosse !!{derivable}, a teoria seria inconsistente. G\"odel
mostrou que a negação dessa !!{sentence} tampouco é !!{derivable} a partir
do sistema de axiomas que ele considerava. (Para esta segunda parte, G\"odel
teve de supor que a teoria~$\Th{T}$ é o que se chama de
``$\omega$-consistente''. A $\omega$-consistência está relacionada à
consistência, mas é uma propriedade mais forte.\footnote{Isto é, toda teoria
$\omega$-consistente é consistente, mas a recíproca não vale.} Alguns anos
depois de G\"odel, Rosser mostrou que basta supor a consistência simples
de~$\Th{T}$.)

O primeiro desafio é compreender como construir !!a{sentence} que se refira
a si mesma. Para toda !!{formula}~$!A$ na linguagem de~$\Th{Q}$, seja
~$\gn{!A}$ o numeral correspondente a~$\Gn{!A}$. Pense no que isso
significa: $!A$~é !!a{formula} na linguagem de~$\Th{Q}$,
$\Gn{!A}$~é um número natural, e $\gn{!A}$~é um \emph{termo} na linguagem
de~$\Th{Q}$. Assim, toda !!{formula}~$!A$ na linguagem de~$\Th{Q}$ tem um
\emph{nome},~$\gn{!A}$, que é um termo na linguagem de~$\Th{Q}$; isso nos
fornece um quadro conceitual no qual !!{formula}s da linguagem de~$\Th{Q}$
podem ``dizer'' coisas sobre outras !!{formula}s. O lema seguinte é
conhecido como lema do ponto fixo.

\begin{lem}
Seja $\Th{T}$ qualquer teoria que estenda~$\Th{Q}$, e seja $!B(x)$ qualquer
!!{formula} cuja única variável livre seja~$x$. Então existe
!!a{sentence}~$!A$ tal que $\Th{T} \Proves!A \liff !B(\gn{!A})$.
\end{lem}

O lema afirma que, dada qualquer propriedade $!B(x)$, existe
!!a{sentence}~$!A$ que afirma ``$!B(x)$ vale para mim'', e $\Th{T}$
``sabe'' disso.

Como podemos construir tal !!a{sentence}? Considere a seguinte versão do
paradoxo de Epimênides, devida a Quine:
\begin{quote}
``Produz falsidade quando precedida por sua citação'' produz falsidade
quando precedida por sua citação.
\end{quote}
Essa sentença não é diretamente autorreferencial. Ela apenas faz uma
afirmação sobre os objetos sintáticos entre aspas e, ao fazê-lo, está em pé
de igualdade com sentenças como
\begin{enumerate}
\item ``Roberto'' é um nome bonito.
\item ``Eu corri.'' é uma sentença curta.
\item ``Tem três palavras'' tem três palavras.
\end{enumerate}
Mas o que acontece quando tomamos a expressão ``produz falsidade quando
precedida por sua citação'' e a precedemos por uma versão entre aspas de si
mesma? Obtemos a sentença original!{} Em suma, a sentença afirma que é
falsa.

\end{document}

